% ---
% id: meat-core/def:dropped
% kind: def
% title: "Dropped set and dropped program"
% summary: >-
%   A run may give up the entropy of chosen pairs of a key-live acceptance
%   value and a joint sector by setting their deficits to zero, keeping every
%   variable, statistic and constraint. The result is a program of its own;
%   once any pair is dropped, its value is never reported as the undropped
%   one.
% prominence: standard
% status: implemented
% uses: [def:joint-sectors, def:feasible, thm:branch-formula, def:declaration]
% mentions: [cor:drop-program, prop:retention, lem:joint-sectors]
% assumes: []
% refs: []
% provenance: newly-derived
% implements: certiqs_key.meat_core.objective.joint:joint_objective_value
% lean:
% notes: "No linear surrogate constant is defined here; the one proved instance is lem:rank-one-sector."
% ---

A run may prefer a weaker bound that is cheaper to compute: a
\tref{joint sector} whose deficits it gives up for every
\tref{key-live acceptance value} needs no matrix function (Corollary~\ref{cor:drop-program}\,(iii)). What the run
gives up is only the entropy those deficits carry; the sector stays in the
program.

\begin{definition}\label{def:dropped}
Let $\Decl$ be a declaration in the general format of
Definition~\ref{def:declaration}. A
\term[A set of pairs of a key-live acceptance value and a joint sector whose entropy deficits a run replaces by zero.]{dropped set}
is a set
$\mathcal{D}\subseteq\Chkey\times(\mathbb{T}\times\SecSet)$ of pairs of a
\tref{key-live acceptance value} and a \tref{joint sector}. For
$\omega=\bigoplus_{(t,s)}\omega^{(t,s)}$ block diagonal along the joint
sectors put
\begin{equation}\label{eq:KD}
  K^{\mathcal{D}}_{\hat c}(\omega)\coloneqq P_{\hat c}(\omega)
  +\sum_{(t,s):\,(\hat c,(t,s))\notin\mathcal{D}}
  f_{\hat c,(t,s)}\big(\omega^{(t,s)}\big)
  \qquad(\hat c\in\Chkey),
\end{equation}
with $f_{\hat c,(t,s)}$ of \eqref{eq:joint-deficit}; let
$O^{\mathcal{D}}_{\alpha}$ be \eqref{eq:objective} with
$K^{\mathcal{D}}_{\hat c}$ in place of $K_{\hat c}$, with the same
conventions at vanishing arguments; and let
\begin{equation}\label{eq:hlowD}
  \underline{h}^{\mathcal{D}}_{\alpha}\coloneqq
  \inf_{\mathbf q\in\acSet}\ \inf_{\omega\in\Feas}\
  O^{\mathcal{D}}_{\alpha}(\mathbf q,\omega).
\end{equation}
The minimisation in \eqref{eq:hlowD} is the
\term[The key-rate program with the deficits of a dropped set replaced by zero, over the unchanged feasible set.]{dropped program}.
A joint sector $(t,s)$ is
\term[Said of a joint sector whose pairs with every key-live acceptance value all lie in the dropped set.]{fully dropped}
when $\Chkey\times\{(t,s)\}\subseteq\mathcal{D}$. For
$\mathcal{D}=\varnothing$ these are $K_{\hat c}$ (\eqref{eq:K-joint}),
$O_{\alpha}$ and $\hlow$.

\emph{What is not changed.} The \tref{feasible set} is $\Feas$, unchanged: a
dropped sector's block remains a variable, enters $\Phimap$ and every
$P_{\hat c}$ with its full weight --- the term
$P_{\hat c,(t,s)}(\omega^{(t,s)})$ stays inside $P_{\hat c}$ in
\eqref{eq:KD} --- and enters the marginal condition and every constraint
row. Setting a dropped block to zero, omitting its statistics, or removing
its weight from $P_{\hat c}$ in \eqref{eq:KD} are different operations, and
each can overstate the value (Proposition~\ref{prop:retention} for the
first). \emph{The dropped set is a formulation choice of a run and not
\tref{declared data}}.
\end{definition}
